\begin{frame}
  \onslide<1->{\frametitle{Damped Harmonic Oscillator}
  \framesubtitle{Spring, Mass, Damper System}}
  \onslide<2->{
\begin{figure}
  \includegraphics[width=0.85\textwidth]{\dayonesegtwo/damped_harmonic_oscillator.png}
\end{figure}
}
\onslide<3->{\center --------------------------------------------------------------------}
\end{frame}


\begin{frame}
  \onslide<1->{\frametitle{Damped Harmonic Oscillator} 
  \framesubtitle{System Dynamics}}
\begin{itemize}
  \onslide<2->{\item System model:  Number of states = 2 (position and velocity) }
  \onslide<3->{Derived from differential equation model: $\damposcdiffeq{x}$}
  \onslide<4->{\[\begin{bmatrix}
x_{1}(j)\\
x_{2}(j)
\end{bmatrix}=
\exp
\left(
\begin{bmatrix}
0 & 1\\
-b/m & -k/m
\end{bmatrix}
 \Delta t
\right)
\begin{bmatrix}
x_{1}(j-1)\\
x_{2}(j-1)
\end{bmatrix}
+
\begin{bmatrix}
w_{1}(j)\\
w_{2}(j)
\end{bmatrix}
\]}
    \onslide<4->{\item Thus, $\stm{\Delta t}=\exp
\left(
\begin{bmatrix}
0 & 1\\
-b/m & -k/m
\end{bmatrix}
 \Delta t
\right)$.   }
  \onslide<5->{\item Note that $\exp$ is a matrix exponential. We'll explain later how to derive state transition 
matrices from the underlying differential equations}
  \onslide<6->{\item Here $\Delta t=t_j - t_{j-1}$, $m$ is the mass, $k$ is the 
spring constant, and $b$ is the damping constant; we are using $t$ for the time step index, to avoid clash with spring constant}
\end{itemize} 
\onslide<7->{\center --------------------------------------------------------------------}
\end{frame}

\begin{comment}
\begin{frame}
  \onslide<1->{\frametitle{Damped Harmonic Oscillator}
  \framesubtitle{Process Noise}}
\begin{itemize}
  \onslide<2->{\item Assume that the process noise $w$ has the following statistics:
\[Q_t=\begin{bmatrix}
q_{11} & 0\\
0 & q_{22}
\end{bmatrix}
\]}
  \onslide<3->{  \item This means that there is some process noise in both position and velocity.}
  \onslide<4->{  \item It also means that the noise sequences in position and velocity are statistically 
uncorrelated (i.e., no non-zero off-diagonal elements in $Q_t$.)}
  \onslide<5->{  \item One may think of process noise as due to unmodelled dynamic effects}
\end{itemize} 
\onslide<6->{\center --------------------------------------------------------------------}
\end{frame}
\end{comment}

\begin{frame}
  \onslide<1->{\frametitle{Dampled Harmonic Oscillator}
   \framesubtitle{Measurement Model}}
\begin{itemize}
  \onslide<2->{  \item Measurement Model:
\[z(j) = \begin{bmatrix}1 & 0\end{bmatrix}\begin{bmatrix}
x_{1}(j)\\
x_{2}(j)
\end{bmatrix}+
v(j)
\]
where the measurement noise $v(j)$ has variance $\sigma^2$}
  \onslide<3->{  \item Thus $H_t=\begin{bmatrix}1 & 0\end{bmatrix}$}
  \onslide<4->{  \item Note that we only have a position measurement (similar to constant velocity example).}
\end{itemize} 
\onslide<5->{\center --------------------------------------------------------------------}
\end{frame}

\begin{frame}
  \onslide<1->{\frametitle{Damped Harmonic Oscillator}
  \framesubtitle{State Transition Matrix $\Phi$}}
\begin{itemize}
  \onslide<2->{  \item $\stm{\Delta t}$ can be approximated through series expansion of 
\[\exp\begin{bmatrix}
0 & \Delta t\\
-\dfrac{b}{m}\Delta t & -\dfrac{k}{m}\Delta t
\end{bmatrix}
\]}
  \onslide<3->{  \item Or we can use Matlab's (or scipy's) expm function for specific values 
  of $k$,$b$,$m$, and $\Delta t$.}
\end{itemize} 
\onslide<4->{\center --------------------------------------------------------------------}
\end{frame}

\begin{frame}
  \onslide<1->{\frametitle{Damped Harmonic Oscillator}
  \framesubtitle{Covariance Extrapolation}}
\begin{itemize}
  \onslide<2->{  \item Because $\Phi$ and $Q$ are matrices, the state and covariance extrapolation equations do not simplify from their general matrix form}
  \onslide<3->{  \item Moreover, like $\Phi$, which is derived from a continuous model, $Q$ is also dependent on $\Delta t$ and $q$, the continuous process noise variance:
\begin{align*}
Q &=
q\begin{bmatrix}
\dfrac{(\Delta t)^3}{3} & \dfrac{(\Delta t)^2}{2}\\
\dfrac{(\Delta t)^2}{2} & \Delta t
\end{bmatrix}
\end{align*}
}
  \onslide<4->{\item The details of how this is derived from the continuous model are beyond the scope of the current presentation.  See, e.g., Grewal and Andrews.}
\end{itemize} 
\onslide<5->{\center --------------------------------------------------------------------}
\end{frame}

\begin{frame}
  \onslide<1->{\frametitle{Damped Harmonic Oscillator}
	\framesubtitle{Covariance Extrapolation}}
	\begin{itemize}
		\onslide<2->{\item $\altcovextrapeqarg{j}$ becomes:  }
		{\tiny \begin{align*} 
		\onslide<3->{
			\covmat{j}{-} &=
			\begin{bmatrix}
			\phi_{11} & \phi_{12}\\
			\phi_{12} & \phi_{22}
			\end{bmatrix} 
			\begin{bmatrix}
			p_{11} & p_{12}\\
			p_{21} & p_{22}
			\end{bmatrix}
			\begin{bmatrix}
			\phi_{11} & \phi_{21}\\
			\phi_{12} & \phi_{22}
			\end{bmatrix}
			+\begin{bmatrix}
			q_{11} & q_{12}\\
			q_{21} & q_{22}
			\end{bmatrix}
		}\\
		\onslide<4->{
			&= 		
\begin{bmatrix}
(\phi_{11}p_{11}+\phi_{12}p_{21})\phi_{11}+(\phi_{11}p_{12}+\phi_{12}p_{22})\phi_{12} & (\phi_{11}p_{11}+\phi_{12}p_{21})\phi_{21}+(\phi_{11}p_{12}+\phi_{12}p_{22})\phi_{22}\\
(\phi_{12}p_{11}+\phi_{22}p_{21})\phi_{11}+(\phi_{12}p_{12}+\phi_{22}p_{22})\phi_{12} & (\phi_{12}p_{11}+\phi_{22}p_{21})\phi_{21}+(\phi_{12}p_{12}+\phi_{22}p_{22})\phi_{22}
\end{bmatrix}
\\
&\;\;\;\;\;+
\begin{bmatrix}
q_{11} & q_{12}\\
q_{21} & q_{22}
\end{bmatrix}		}\\
		\onslide<5->{
			&=		
\begin{bmatrix}
\phi_{12}p_{12}\phi_{12}+q_{11} & \phi_{12}p_{12}\phi_{21}+\phi_{12}p_{22}\phi_{22}+q_{12}\\
\phi_{22}p_{21}\phi_{12}+q_{21} & \phi_{22}p_{21}\phi_{21}+\phi_{22}p_{22}\phi_{22}+q_{22}
\end{bmatrix}
		}
		\end{align*}} 
		\onslide<6->{\item Note that the last matrix results when we simplify using $\phi_{11}=0$ (See previous slide)  }
		\onslide<7->{\item We've also decluttered by suppressing the + superscripts and $j-1$ indices on the right-hand sides}
	\end{itemize}
	\onslide<8->{\center --------------------------------------------------------------------}
\end{frame}



\begin{frame}
	\onslide<1->{\frametitle{Damped Harmonic Oscillator}
		\framesubtitle{Kalman Gain}}
	\begin{itemize}
		\onslide<2->{  \item The gain equation $\genkalgaineq$ looks like this:}
			\begin{align*}
			\onslide<3->{ K &=
			\begin{bmatrix}
			p_{11} & p_{12}\\
			p_{21} & p_{22}
			\end{bmatrix}
			\begin{bmatrix}1\\
			0
			\end{bmatrix}
			\left( \begin{bmatrix}1 & 0\end{bmatrix}\begin{bmatrix}
			p_{11} & p_{12}\\
			p_{21} & p_{22}
			\end{bmatrix}
			\begin{bmatrix}
			1\\ 
			0
			\end{bmatrix}
			+\sigma^2
			\right)^{-1}\\}
			\onslide<4->{&=
			\begin{bmatrix}
			\dfrac{p_{11}}{p_{11}+\sigma^2}\\
			\dfrac{p_{21}}{p_{11}+\sigma^2}
			\end{bmatrix},}
			\end{align*}
			\onslide<5->{where we have once again suppressed the minus(-) superscripts and time step variable $j$ on the right hand side}
	\end{itemize} 
	\onslide<6->{\center --------------------------------------------------------------------}
\end{frame}


\begin{comment}
\begin{frame}
	\onslide<1->{\frametitle{Damped Harmonic Oscillator}
		\framesubtitle{Kalman Gain}}
	\begin{itemize}
		\onslide<2->{  \item Because $\Phi$ and $Q$ are matrices, the state and covariance extrapolation equations do not simplify from their general matrix form}
		\onslide<3->{  \item However, the gain equation $\genkalgaineq$ looks like this:
			\begin{align*}
			K &=
			\begin{bmatrix}
			p_{11} & p_{12}\\
			p_{21} & p_{22}
			\end{bmatrix}
			\begin{bmatrix}1\\
			0
			\end{bmatrix}
			\left( \begin{bmatrix}1 & 0\end{bmatrix}\begin{bmatrix}
			p_{11} & p_{12}\\
			p_{21} & p_{22}
			\end{bmatrix}
			\begin{bmatrix}
			1\\ 
			0
			\end{bmatrix}
			+\sigma^2
			\right)^{-1}\\
			&=
			\begin{bmatrix}
			\dfrac{p_{11}}{p_{11}+\sigma^2}\\
			\dfrac{p_{21}}{p_{11}+\sigma^2}
			\end{bmatrix}
			\end{figure}
			\end{align*}
			where we have once again suppressed the +/- superscripts and time step variable.}
	\end{itemize} 
	\onslide<4->{\center --------------------------------------------------------------------}
\end{frame}
\end{comment}

\begin{frame}
	\onslide<1->{\frametitle{Damped Harmonic Oscillator Estimator}
		\framesubtitle{MATLAB/Python Run}}
	\begin{itemize}
	\onslide<2->{\item $\Delta t=0.1$ sec, 50 samples,\\ 50\% process noise assumption factor
	}
	\onslide<3->{\item $m=1$ kg 
}
	\onslide<4->{\item Light spring:  $k=1$ N/m
}
	\onslide<5->{\item Medium damping: $b=5$ Ns/m
}
    \end{itemize}
	\onslide<6->{\center --------------------------------------------------------------------}
\end{frame}



\begin{frame}
  \onslide<1->{\frametitle{Damped Harmonic Oscillator Estimator}
  \framesubtitle{MATLAB/Python Run}}
  \onslide<2->{
\begin{figure}
  \includegraphics[width=0.95\textwidth,height=7cm]{\dayonesegtwo/damposc.png}
\end{figure}
}
\onslide<3->{\center --------------------------------------------------------------------}
\end{frame}

\begin{frame}
  \onslide<1->{\frametitle{Damped Harmonic Oscillator Estimator }
  \framesubtitle{Sawtooth Plots}}
  \onslide<2->{
\begin{figure}
  \includegraphics[width=0.95\textwidth,height=7cm]{\dayonesegtwo/damposc_variance.png}
\end{figure}
}
\onslide<3->{\center --------------------------------------------------------------------}
\end{frame}

\begin{frame}
	\onslide<1->{\frametitle{Damped Harmonic Oscillator Estimator}
		\framesubtitle{MATLAB/Python Run}}
	\begin{itemize}
		\onslide<2->{\item $\Delta t=1$ sec, 50 samples,\\ 50\% process noise assumption factor
		}
		\onslide<3->{\item $m=1$ kg 
		}
		\onslide<4->{\item Light spring:  $k=1$ N/m
		}
		\onslide<5->{\item Medium damping: $b=5$  Ns/m
		}
	\end{itemize}
	\onslide<6->{\center --------------------------------------------------------------------}
\end{frame}



\begin{frame}
	\onslide<1->{\frametitle{Damped Harmonic Oscillator Estimator}
		\framesubtitle{MATLAB/Python Run}}
	\onslide<2->{
		\begin{figure}
			\includegraphics[width=0.95\textwidth,height=7cm]{\dayonesegtwo/damposc_dt_1s.png}
		\end{figure}
	}
	\onslide<3->{\center --------------------------------------------------------------------}
\end{frame}

\begin{frame}
	\onslide<1->{\frametitle{Damped Harmonic Oscillator Estimator }
		\framesubtitle{Sawtooth Plots}}
	\onslide<2->{
		\begin{figure}
			\includegraphics[width=0.95\textwidth,height=7cm]{\dayonesegtwo/damposc_variance_dt_1s.png}
		\end{figure}
	}
	\onslide<3->{\center --------------------------------------------------------------------}
\end{frame}

\begin{comment}

\begin{frame}
  \onslide<1->{\frametitle{Damped Harmonic Oscillator Estimator}
  \framesubtitle{MATLAB/Python Run: Conservative Process Noise Estimate}}
  \onslide<1->{
\begin{figure}
  \includegraphics[width=0.95\textwidth,height=7cm]{\dayonesegtwo/damposc_conservative.png}
\end{figure}
}
\onslide<4->{\center --------------------------------------------------------------------}
\end{frame}

\begin{frame}
  \onslide<1->{\frametitle{Damped Harmonic Oscillator Estimator Sawtooth}
  \framesubtitle{MATLAB/Python Run: Conservative Process Noise Estimate}}
  \onslide<1->{
\begin{figure}
  \includegraphics[width=0.95\textwidth,height=7cm]{\dayonesegtwo/damposc_variance_conservative.png}
\end{figure}
}
\onslide<2->{\center --------------------------------------------------------------------}
\end{frame}

\end{comment}